\cotefile{123}
\scanpage{20}
nier (ou, ce qui revient au même, $c$) est bien déterminée. Les unipotents qui définissent ladite filtration, en termes d'un $i$ qui la définit, sont \inserted{ceux} de la forme $\exp(e)$, où $e$ parcourt $\handwritten{\mathfrak{g}}(2)^{\handwritten{*}}+\operatorname{Fil}^{\handwritten{3}}(\handwritten{\mathfrak{g}})$. Ils forment une classe de conjugaison d'éléments de $U$ \uline{par} $P$ (NB mais pas par $U$), les classes de conjugaison \uline{par} $U$ s'identifiant aux ensembles de la forme $(\exp(e)\handwritten{\mid}e\handwritten{\in}a+\operatorname{Fil}^{\handwritten{3}}(\handwritten{\mathfrak{g}},))$, où $a\handwritten{\in\mathfrak{g}}(2)^{\handwritten{*}}$ est un élément quelconque.\deleted{de}

\uline{Cas principal}. Pour qu'on ait\hspace{1.5em}$(1)=0$, il faut et suffit que \deleted{l'homomor} la classe d'homomorphismes $f:\operatorname{Sl}(2,k)\hspace{1.5em}G$ envisagée soit ``\unsure{princi\deleted{a}p\deleted{p}ale|princi\deleted{p}a\deleted{a}ple}'' au sens de Dynkin-Kostant. Alors le parabolique associé $P$ est un Borel (la réciproque n'étant pas vraie \deleted{sans doute}, cf table de Dynkin), et le groupe $Z_o$ est le groupe unité (i.e. $\operatorname{Dyn}_c$ est un torseur sous $G$), \deleted{\unsure{\$|S}}la réciproque n'étant pas vraie sans doute.
\margin{\handwritten{\uline{ex}\quad Cas $A_2,A_4$ les paraboliques associés aux unipotents ``sous-réguliers'' sont des \unsure{Borels|borels} !}}
% archive sheet number: 19
